% ============================================================
% §5.3 Q3结果分析 — 服务定价与政府补贴 MILP (v2: F/H/J配置)
% 竞赛: 电工杯 B题 (2026)
% ============================================================

\section{Q3: 服务定价与政府补贴优化结果}

\subsection{最优定价方案}

基于§3.2构造的S3价格满意度Big-M线性化模型，在固化Q2选址方案（F大型+H中型+J中型，10/10全覆盖）后，以词典序目标（优先最大化加权$S_3$满意度、次优先$\varepsilon=0.0001$弱激励营收最大化）求解MILP。CBC在1秒内收敛至全局最优（32变量、63约束），全部6项服务均实现$S_3=1.000$（平价区间），仅上门护理一项需从基准价30元降至12.64元（-57.9\%）以满足F站的8\%利润率紧约束。最优定价方案见表\ref{tab:q3_pricing}。

\begin{table}[htbp]
\centering
\caption{Q3最优定价方案}
\label{tab:q3_pricing}
\begin{tabular}{lccccc}
\hline
\textbf{服务项目} & \textbf{基准价(元)} & \textbf{最优定价(元)} & \textbf{溢价率(\%)} & \textbf{S3} & \textbf{价格区间} \\
\hline
助餐 & 10 & 10.00 & 0.0 & 1.000 & 平价 \\
日间照料 & 20 & 20.00 & 0.0 & 1.000 & 平价 \\
上门护理 & 30 & \textbf{12.64} & \textbf{-57.9} & 1.000 & 平价 \\
康复理疗 & 28 & 28.00 & 0.0 & 1.000 & 平价 \\
助浴 & 25 & 25.00 & 0.0 & 1.000 & 平价 \\
紧急救助 & 0 & 0.00 & 公益免费 & 1.000 & 公益免费 \\
\hline
\end{tabular}
\end{table}

\subsection{站点利润率约束的绑定诊断}

表\ref{tab:q3_stations}给出了最优定价下的站点级运营指标。F站（大型，日容量3,000人次）的利润率恰好触达8.0\%的监管红线——这是全系统唯一的紧约束（binding constraint）。H站与J站（均为中型，日容量2,000人次）的利润率仅为7.3\%，说明在S3最大化的词典序目标下，两站尚存0.7个百分点的定价空间（即可通过提高约1--2项服务的价格来增加营收），但因$S_3$已全域触及1.000（平价段上线），任何提价行为都会将对应服务的$S_3$从1.000降至0.900（微溢价段），造成加权满意度目标函数的净损失。这一结果完美诠释了词典序目标的设计意图：\textbf{在满足所有站点可持续运营（利润率$\ge 0$且$\le 8\%$）的前提下，S3优先于营收}。

\begin{table}[htbp]
\centering
\caption{最优定价下的站点运营指标}
\label{tab:q3_stations}
\begin{tabular}{lcccccc}
\hline
\textbf{站点} & \textbf{规模} & \textbf{年营收(万元)} & \textbf{年补贴(万元)} & \textbf{年总成本(万元)} & \textbf{年利润(万元)} & \textbf{利润率(\%)} \\
\hline
F & 大型 & 1,967.1 & 93.6 & 1,908.1 & 152.6 & \textbf{8.0 (紧约束)} \\
H & 中型 & 1,265.6 & 64.8 & 1,239.6 & 90.8 & 7.3 \\
J & 中型 & 1,226.3 & 64.8 & 1,203.2 & 87.9 & 7.3 \\
\hline
\textbf{合计} & -- & 4,459.0 & 223.2 & 4,350.9 & 331.3 & 7.6 (加权) \\
\hline
\end{tabular}
\end{table}

\subsection{三层交叉补贴机制分析}

\subsubsection{第一层：站内营利养公益}

Q3的全系统紧急救助月需求为4,940次（F:2,202 + H:1,399 + J:1,339），单价0元（公益免费），单次直接成本8元。全系统紧急救助年净支出为47.42万元——这笔纯公益支出由各站营利性服务的利润（年总利润331.3万元）完全吸收，交叉补贴率仅14.3\%（47.42/331.3），规模经济效应（全覆盖带来的总利润较B/C/J配置增长145\%）显著摊薄了公益服务的单位吸收成本。

\subsubsection{第二层：站间补贴帽的效率错配}

政府补贴规则为"非紧急救助服务每单补贴2元，受站点日补贴上限约束"。在三站满负荷运营的条件下，各站的理论日补贴额与实际上限的对比见表\ref{tab:q3_subsidy}。三个站点的有效补贴率分别为0.73元/次（F）、0.79元/次（H）、0.81元/次（J），均远低于理论的2元/次——F站因规模最大（需求最多），效率损失最为严重（63.5\%的潜在补贴被上限截断），形成一种"逆向再分配"：大型站因服务人次多而触发补贴帽更早，有效补贴率反而最低。

\begin{table}[htbp]
\centering
\caption{各站点补贴池约束分析}
\label{tab:q3_subsidy}
\begin{tabular}{lccccc}
\hline
\textbf{站点} & \textbf{日非紧急需求(次)} & \textbf{理论补贴(元/日)} & \textbf{日补贴上限(元)} & \textbf{实际补贴(元/日)} & \textbf{有效补贴率(元/次)} \\
\hline
F & 3,559 & 7,118 & 2,600 & 2,600 & 0.73 \\
H & 2,291 & 4,582 & 1,800 & 1,800 & 0.79 \\
J & 2,221 & 4,442 & 1,800 & 1,800 & 0.81 \\
\hline
\end{tabular}
\end{table}

\subsubsection{第三层：上门护理的"拉格朗日影子价格"效应}

在Q3的最优解中，上门护理（30元$\to$12.64元，-57.9\%）成为唯一的降价服务，而助餐（频次最高、需求弹性最大）维持基准价不动。这一结果与F/H/J的站点规模结构直接相关：F站（大型）是利润率紧约束的绑定者，而上门护理在所有服务中具有最高的基准价（30元），其单位降价的利润率释放效应最大——从拉格朗日乘子的视角，上门护理的"影子价格"（降低1元对利润率松绑的边际贡献）在所有服务中最高，因此成为词典序S3最大化目标下的最优"价格牺牲品"。

从Ramsey-Boiteux最优定价的视角看：在总利润率约束下，价格需求弹性最小的服务应承担最多的加价（或最少的降价），反之亦然。上门护理的基准价（30元）在所有服务中最高，其"过剩的利润空间"使其成为$S_3$最大化目标下最优先被削减的对象。这一结果可理解为"反向价格歧视"——高频低价服务（助餐：10元$\times$118,443次/月）的保护优先于低频高价服务（上门护理：30元$\times$19,608次/月）的保护，在社会福利函数中体现了"利用最不敏感价格维度吸收监管约束"的Ramsey原理。

\subsection{加权综合满意度解构}

表\ref{tab:q3_satisfaction}将各小区的三维满意度分解为$S_1$（距离，已由Q2分配固定）、$S_2$（响应满意度，全部0.600因三站满负荷）、$S_3$（价格满意度，全部1.000因Q3最优定价维持平价）。加权综合满意度$S$的跨小区变异系数仅为0.035（区间[0.800, 0.880]），反映了全覆盖方案下满意度的空间均等化特征。

\begin{table}[htbp]
\centering
\caption{各小区加权综合满意度分解（Q2+Q3联合）}
\label{tab:q3_satisfaction}
\begin{tabular}{lcccccc}
\hline
\textbf{小区} & \textbf{服务站} & $S_1$(距离) & $S_2$(响应) & $S_3$(价格) & $S$(综合) & \textbf{说明} \\
\hline
A & J & 0.900 & 0.600 & 1.000 & 0.860 & 中距离，满意均好 \\
B & H & 0.900 & 0.600 & 1.000 & 0.860 & 中距离，满意均好 \\
C & F & 0.600 & 0.600 & 1.000 & 0.800 & 远距离S1触底 \\
D & J & 0.900 & 0.600 & 1.000 & 0.860 & 中距离，满意均好 \\
E & H & 0.600 & 0.600 & 1.000 & 0.800 & 边界距离S1触底 \\
F & F & 1.000 & 0.600 & 1.000 & 0.880 & 本小区自服务 \\
G & F & 0.750 & 0.600 & 1.000 & 0.830 & 中距离，满意均好 \\
H & H & 1.000 & 0.600 & 1.000 & 0.880 & 本小区自服务 \\
I & F & 0.600 & 0.600 & 1.000 & 0.800 & 远距离S1触底 \\
J & J & 1.000 & 0.600 & 1.000 & 0.880 & 本小区自服务 \\
\hline
\textbf{平均} & -- & 0.825 & 0.600 & 1.000 & 0.845 & -- \\
\hline
\end{tabular}
\end{table}
